\section{The prompts and the output contract}
\label{sec:prompts}

This paper measures what a model adds to a pipeline, and that measurement cannot be read without
knowing what the model was asked. The full text of the four prompts is in the released repository
\cite{release2026artefact}; what is reproduced here is the part the paper's own claims rest on,
because a claim about falsification before confidence is a claim about specific sentences in a
specific prompt.

\subsection{What the three analyst prompts share}\label{prompt-shape}

Each analyst receives one system prompt with the same five parts, and the parts are the same
length in each: a role and an evidence rule, a tool workflow, a set of conditional tools, a
verification discipline, and a set of standing cautions. The evidence rule is identical across
the three and is the pipeline's grounding constraint: for every claim the model makes it must
cite a concrete artefact, which for the static analyst is a function name, a string offset, an
import or a hex pattern, for the dynamic analyst an API call or a registry write from the sandbox
report, and for the network analyst a flow or a resolved domain. What differs between the three
is the tool server each is bound to, the evidence type it reads, and the technique families its
role paragraph names.

\subsection{The confidence rules, as shipped}\label{prompt-confidence}

Two clauses of the verification discipline carry the fourth architectural claim of
Section~\ref{sec:contributions} and set the cap whose firing rate Section~\ref{sec:mechanisms}
measures at \fact{cap-rate} of techniques. They are quoted from the static analyst's prompt
without alteration:

\begin{verbatim}
- A SPECIFIC claim (a named algorithm like RC4/djb2/ROR13, a constant
  or XOR key, or a hash-resolved API) may reach CONFIDENCE >= 0.8 only
  if you FALSIFY it first: `emulate_function` with a known input vs the
  expected output, OR `analyze_dataflow(direction=backward)` to confirm
  its origin. If you cannot run the check (non-leaf, syscall/heap side
  effects), cap CONFIDENCE at 0.7.

- A claim is High (>= 0.8) only with >= 2 independent evidence loci
  (e.g. an import AND its call-site). A single locus caps at 0.7.
  Reconcile any contradictory signals before emitting.
\end{verbatim}

Two further clauses of the same block are elided rather than silently dropped: one governs
disambiguation when a hash-resolution tool returns colliding candidates, and one refuses a
particular obfuscation technique on the evidence of dynamic import resolution alone. Neither
bears on a measurement reported here, and both are in the release.

The rules are worth reading against the result. They ask the model to earn a number above
\setting{0.8} by running a check, and to cap itself at \setting{0.7} when it cannot. The number
they produce is the one every deterministic gate downstream consumes, and
Section~\ref{sec:mechanisms} finds it near chance, with \fact{confidence-modal-count} of
\fact{confidence-n} claims arriving at exactly \setting{1.0}. Whether the discipline is not
followed, or is followed and does not separate correct claims from wrong ones, this design cannot
distinguish.

\subsection{The output contract}\label{prompt-schema}

An analyst answers in blocks, and a block is the unit everything downstream counts. The parser
requires the first field and defaults the rest, because a block that names a finding and omits
its confidence is still a finding:

\begin{verbatim}
CLAIM: <the finding>
EVIDENCE: <a concrete artefact reference>
CONFIDENCE: <0.0 to 1.0>
TECHNIQUE: <T1055.001 | NONE>
---
\end{verbatim}

Each parsed block becomes one record, and its four fields are what the rest of the system sees.
The claim and the evidence reference are free text. The confidence is bounded to the unit
interval by the schema rather than by the prompt, which is why a model that ignores the
verification discipline still produces a well-formed number. The technique identifier is optional
and is pattern-checked where it is declared:

\begin{verbatim}
technique_id: str | None = Field(None, pattern=r"^T\d{4}(\.\d{3})?$")
\end{verbatim}

A malformed identifier is therefore refused at the schema and never reaches the cascade. This is
the filter that Section~\ref{sec:cascade} finds the text-fallback path bypassing, which is why
the identifiers admitted there are real ones no evidence source claimed rather than invented
strings.

\subsection{What varies per sample}\label{prompt-interpolation}

The system prompts are fixed strings and are not templated per sample, which is the reason
Section~\ref{sec:limitations} records prompt sensitivity as \texttt{PARTIAL}: production prompts
were held constant and were not varied per model. What varies is the human turn, which carries
the sample's path for the static analyst, the sandbox report for the dynamic analyst and the
network observations for the network analyst, each wrapped in delimiters that mark it as
untrusted input. Evidence is truncated to a per-call character cap before it is interpolated,
and that cap is one of the truncation points enumerated under P6.
